\chapter*{Bảng thuật ngữ và từ viết tắt}
\addcontentsline{toc}{chapter}{Bảng thuật ngữ và từ viết tắt}

Khóa luận giữ nguyên một số thuật ngữ chuyên ngành bằng tiếng Anh khi bản
dịch tiếng Việt chưa phổ biến hoặc dễ gây hiểu nhầm. Bảng dưới đây chú thích
ý nghĩa của các thuật ngữ đó (xếp theo thứ tự bảng chữ cái), tiếp theo là
bảng các từ viết tắt.

\section*{Bảng chú thích thuật ngữ}

\begin{longtable}{L{4.2cm} L{9.8cm}}
\toprule
\textbf{Thuật ngữ} & \textbf{Giải nghĩa} \\
\midrule
accessibility tree (axtree) & Cây trợ năng: biểu diễn dạng văn bản của các
phần tử giao diện đang hiển thị trên một trang web, mỗi phần tử kèm một mã
định danh; là ``quan sát'' mà agent nhìn thấy thay cho ảnh chụp màn hình. \\
action space & Không gian hành động: tập hữu hạn các thao tác agent được
phép thực hiện (click, điền form, điều hướng, \dots). \\
agent & Tác nhân phần mềm tự trị: chương trình tự quan sát môi trường,
suy luận và hành động để đạt một mục tiêu. \\
benchmark & Bộ chuẩn đánh giá: môi trường + tập task + thang đo thống nhất
để so sánh các phương pháp. \\
best practice & Kinh nghiệm thực hành tốt đã được kiểm chứng, nên áp dụng
lại trong các dự án tương tự. \\
chat template & Khuôn dạng hội thoại mà một mô hình ngôn ngữ được huấn luyện
để đọc (cách đánh dấu vai trò, lượt nói, chế độ suy nghĩ). \\
crawl / crawling & Duyệt và thu thập trang web tự động theo liên kết. \\
curriculum & Trình tự huấn luyện sắp xếp dữ liệu từ dễ đến khó. \\
fine-tune / fine-tuning & Tinh chỉnh: huấn luyện tiếp một mô hình đã
pretrain trên dữ liệu chuyên biệt để chuyên hóa nó cho một nhiệm vụ. \\
frontier & Biên khám phá: hàng đợi các URL ứng viên chưa thăm, quyết định
``đi đâu tiếp theo'' trong quá trình khám phá website. \\
judge & Bộ chấm độc lập: một LLM thứ hai kiểm tra lại các episode mà agent
tự tuyên bố thành công. \\
low-level task & Task mức thấp: nhiệm vụ ngắn, gắn với một trang cụ thể, do
LLM đề xuất khi khám phá (phân biệt với semantic task). \\
operator & Toán tử lập kế hoạch: một skill được nâng lên dạng
$\langle$precondition, effect$\rangle$ (điều kiện tiên quyết, hiệu ứng) để
có thể ghép chuỗi bằng thuật toán lập kế hoạch. \\
option & Khung hình thức hóa skill trong học tăng cường: bộ ba $\langle I,
\pi, \beta \rangle$ gồm tập khởi đầu, chính sách và điều kiện kết thúc. \\
pipeline & Chuỗi xử lý: dãy các bước tự động nối tiếp nhau, đầu ra của bước
trước là đầu vào của bước sau. \\
precondition / effect & Điều kiện tiên quyết / hiệu ứng của một operator,
biểu diễn bằng các mệnh đề ngôn ngữ tự nhiên đã được chuẩn hóa. \\
prompt & Toàn bộ văn bản đầu vào đưa cho mô hình ngôn ngữ trong một lần
gọi (bao gồm chỉ dẫn hệ thống, quan sát, lịch sử). \\
prompt parity & Nguyên tắc ``đồng nhất prompt'': prompt trong dữ liệu huấn
luyện phải trùng khớp từng cấu trúc với prompt agent nhìn thấy lúc chạy
thật. \\
self-host / self-hosted & Tự triển khai trên hạ tầng của mình thay vì dùng
dịch vụ bên thứ ba. \\
semantic task & Task ngữ nghĩa: nhiệm vụ dài hạn được sinh ra từ operator
graph bằng cách ghép chuỗi. \\
skill & Kỹ năng: một hành vi tái sử dụng được, có tham số, được phát hiện
từ các trajectory lặp lại và kiểm chứng trước khi đưa vào thư viện. \\
task & Nhiệm vụ giao cho agent, phát biểu bằng ngôn ngữ tự nhiên. \\
trajectory & Quỹ đạo tương tác: chuỗi các cặp (quan sát, hành động) sinh ra
khi agent thực hiện một task. \\
web agent & Tác nhân web: agent thao tác website qua giao diện người dùng
(đọc trang, click, gõ phím) như một người dùng thật. \\
workflow & Chuỗi thao tác nghiệp vụ lặp lại theo một trình tự quen thuộc. \\
world model & Mô hình thế giới: mô hình học cách dự đoán quan sát kế tiếp
của môi trường, dùng để sinh dữ liệu mô phỏng. \\
wrap-up call & Lời gọi LLM ``chốt hạ'' cuối episode khi hết ngân sách bước,
chỉ được phép trả lời hoàn thành hoặc bất khả thi. \\
\bottomrule
\end{longtable}

\section*{Bảng từ viết tắt}

\begin{longtable}{L{3.5cm} L{10.5cm}}
\toprule
\textbf{Viết tắt} & \textbf{Nghĩa đầy đủ} \\
\midrule
API & Application Programming Interface - giao diện lập trình ứng dụng \\
AWM & Agent Workflow Memory \\
BFS & Breadth-First Search - tìm kiếm theo chiều rộng \\
CI & Confidence Interval - khoảng tin cậy \\
CoT & Chain-of-Thought - chuỗi suy luận \\
DFS & Depth-First Search - tìm kiếm theo chiều sâu \\
DOM & Document Object Model - mô hình đối tượng tài liệu \\
HSGE & Hierarchical Skill-Graph Exploration - khám phá dựa trên graph kỹ
năng phân tầng (phương pháp của khóa luận) \\
LLM & Large Language Model - mô hình ngôn ngữ lớn \\
LoRA & Low-Rank Adaptation - tinh chỉnh hạng thấp \\
MDP & Markov Decision Process - quá trình quyết định Markov \\
NL & Natural Language - ngôn ngữ tự nhiên \\
POMDP & Partially Observable MDP - MDP quan sát được một phần \\
RL & Reinforcement Learning - học tăng cường \\
SFT & Supervised Fine-Tuning - tinh chỉnh có giám sát \\
SLM & Small Language Model - mô hình ngôn ngữ nhỏ \\
SoM & Set-of-Marks - kỹ thuật đánh dấu phần tử trên ảnh chụp màn hình \\
URL & Uniform Resource Locator - địa chỉ tài nguyên web \\
VLM & Vision-Language Model - mô hình thị giác-ngôn ngữ \\
\bottomrule
\end{longtable}
